% Q16_gold_answers.tex — GRADER CALIBRATION FIXTURE (correct, paraphrased).
% A drop-in replacement for paper/Q16_answers.tex (spec §4.7, §6.12): every answer is
% correct, but the lemma names, the wording and the order of the explanation are
% deliberately far from the reference, so the grader must recognise correctness by
% substance (the rubric's probes), not by matching text.
%
% EXPECTED[B1][statement-fidelity][4], [induction-soundness][4], [case-exhaustiveness][4], [blueprint-quality][4], [gap-accounting][4]: correct; every invented statement is true (each is the statement of a lemma the reference twins prove under another name), every invariant is stated with its bounds and its three duties, and the instance and the named facts compose. Probe A on H1's generalization at b = [], acc = [1], k = 1 fails at its postcondition (no refutation; captured 2026-10-11, Dafny 4.11.0). No finding.
% EXPECTED[B2][card-fidelity][4]: every field agrees with the pinned range (lines 458-481); Variance correctly says what the range does not give. No finding on any field.
% Expected per-problem grades: B1 = 4, B2 = 4 (3 is tolerated only with a reading note and no flaw).
% A calibration run FAILS if any problem scores <= 2 or if any finding claims a false fact,
% a false statement, a wrong instrument outcome or a contradiction with the range.
\documentclass[11pt]{article}
\input{preamble}

\title{Q16 --- Practice III-A \\ \large Written answers}
\author{learner}
\date{}

\begin{document}
\maketitle
\sloppy

% ---------------------------------------------------------------------------
\section*{Problem B1 (informal proof first: the blueprint)}\label{ans:B1}
\begin{answer}
\textbf{H1 (\dfy{MulCorrect}).}
(a) The plain statement cannot be its own hypothesis: unfolding
\dfy{Times(a, b)} once leaves \dfy{MulAcc(a, b[1..], row, 1)}, with a row in
the accumulator and the position at 1, while the hypothesis at \dfy{b[1..]}
is about \dfy{MulAcc(a, b[1..], [], 0)}. So I prove, for every valid
\dfy{acc} and every \dfy{k}, by induction on $|b|$:
\[ \mathit{Valid}(\mathit{MulAcc}(a, b, acc, k)) \;\wedge\;
   \mathit{ToNat}(\mathit{MulAcc}(a, b, acc, k)) = \mathit{ToNat}(acc) +
   \mathit{Pow}(k)\,(\mathit{ToNat}(a)\cdot\mathit{ToNat}(b)). \]
(b) Lemmas: \dfy{AccLaw(a, b, acc, k)}, the statement above
(\dfy{requires Valid(a) && Valid(b) && Valid(acc)}, \dfy{decreases |b|});
\dfy{ShiftLaw(s, k)}: \dfy{requires Valid(s)}, \dfy{ensures Valid(Shift(s, k))
&& ToNat(Shift(s, k)) == Pow(k) * ToNat(s)}, by induction on $k$, using
\dfy{Shift(s, k) == [0] + Shift(s, k - 1)} and L2's \dfy{ToNatCons},
\dfy{ValidCons}.
(c) \dfy{ShiftLaw} first. \dfy{AccLaw}: if $b$ is empty, \dfy{MulAcc} returns
\dfy{acc} and both products are 0 (L4's \dfy{MulZero}). Otherwise, with $A =
\mathit{ToNat}(a)$, $B' = \mathit{ToNat}(b[1..])$, $P = \mathit{Pow}(k)$: the row
\dfy{Shift(MulLimb(a, b[0]), k)} is valid and worth $P(A b_0)$ (L4's
\dfy{MulLimbCorrect}, then \dfy{ShiftLaw}); adding it keeps validity and adds
the values (L4's \dfy{AddCorrect}); the hypothesis at \dfy{b[1..]},
\dfy{Add(acc, row)}, \dfy{k + 1} gives the rest, and $\mathit{Pow}(k + 1) = B P$
by one unfolding. \dfy{MulCorrect} is \dfy{AccLaw} at \dfy{acc = []},
\dfy{k = 0}.
(d) $\mathit{ToNat}(b) = b_0 + \mathrm{BASE}\cdot B'$ (L2's
\dfy{ToNatHeadTail}, \dfy{ValidTail} for \dfy{b[1..]}); the step's identity
$P(A b_0) + (\mathrm{BASE}\,P)(A B') = P(A (b_0 + \mathrm{BASE}\,B'))$ from
\dfy{LemmaMulIsDistributive}, \dfy{LemmaMulIsAssociative} and
\dfy{LemmaMulIsCommutative}; \dfy{MulNonneg} for the \dfy{nat} products;
\dfy{LemmaMulBasics} for $1 \cdot x$ at the instance.
(e) No gap.

\textbf{H2 (\vrs{div_limb}).}
(a) First a vector \vrs{q} of $n = |a|$ zeros. Then $i$ from $n$ down to 0,
with invariant $i \le n$, $|q| = n$, $\mathit{rem} < d$ and
$\mathit{val}(a[i..n]) = \mathit{val}(q[i..n]) \cdot d + \mathit{rem}$ (both
suffixes empty and $\mathit{rem} = 0$ on entry; at $i = 0$ it is the
\vrs{ensures}, after \vrs{normalize}).
(b) \vrs{suffix_peel(s, i)}: for $0 \le i < |s|$,
$\mathit{val}(s[i..]) = s_i + \mathrm{BASE}\cdot\mathit{val}(s[i + 1..])$,
proved by naming \vrs{s[i..].drop_first() =~= s[i + 1..]}.
\vrs{digit_ok(r, x, d)}: for $0 < d < \mathrm{BASE}$, $r < d$,
$x < \mathrm{BASE}$ and $t = r\,\mathrm{BASE} + x$: $t / d < \mathrm{BASE}$,
$t \bmod d < d$, $t = (t / d)\,d + t \bmod d$.
(c) Each step: $i$ goes down by one; \vrs{digit_ok} at the remainder and the
limb, so the 128-bit $t$ splits into a quotient limb $q_i$ and a remainder
$r' < d$; after writing $q_i$, \vrs{suffix_peel} on $a$ and on $q$ (whose part
above $i$ did not change) turns the old equation into
$a_i + \mathrm{BASE}(Q' d + \mathit{rem}) = (q_i + \mathrm{BASE}\,Q') d + r'$.
(d) \vrs{lemma_fundamental_div_mod}, \vrs{lemma_mod_pos_bound} and
\vrs{lemma_multiply_divide_lt} from \vrs{vstd::arithmetic::div_mod}; one scoped
\vrs{nonlinear_arith} step for $t < d\,\mathrm{BASE}$; the step's equation from
\vrs{lemma_mul_is_distributive_add}, its \vrs{_other_way} form and
\vrs{lemma_mul_is_associative}.
(e) No gap.

\textbf{H3 (\vrs{mul} and \dfy{MulLoop}).}
(a) \vrs{mul}: a vector of $n = |a|$ zeros (value 0 by L1's
\vrs{lemma_seq_to_nat_push_zero} at each push). Outer loop over $j < m$:
$|res| = n + j$ and $\mathit{val}(res) = \mathit{val}(a)\cdot\mathit{val}(b[0..j])$.
Inner loop over $i < n$, with $v_0$ the value when the row starts and $c$
the carry: $|res| = n + j$, the outer facts restated (Verus loops do not keep
them), and $\mathit{val}(res) + W(j + i)\,c = v_0 +
W(j)\,(\mathit{val}(a[0..i])\, b_j)$. After the row the carry is pushed.
\dfy{MulLoop}: $0 \le k \le |b|$ and \dfy{MulAcc(a, b[k..], r, k) == Times(a, b)}.
(b) \vrs{value_push(s, x)}: $\mathit{val}(s.push(x)) = \mathit{val}(s) + W(|s|)\,x$
(from L4's \vrs{lemma_seq_to_nat_append}). \vrs{value_write(s, p, x)}: for
$0 \le p < |s|$, $\mathit{val}(s[p := x]) + W(p)\,s_p = \mathit{val}(s) + W(p)\,x$,
by induction on $p$.
(c) Inner step: $a_i b_j + res_{i+j} + c \le \mathrm{BASE}^2 - 1$
(\vrs{lemma_mul_upper_bound}), so the 128-bit sum gives a low limb written at
$i + j$ and a carry below \texttt{BASE}; \vrs{value_write} at $i + j$, the
prefix of $a$ grown by \vrs{value_push}, $W(j + i + 1) = \mathrm{BASE}\,W(j + i)$
and $W(j + i) = W(j)\,W(i)$ (\vrs{lemma_weight_succ}, \vrs{lemma_weight_add}),
then the product identities by name. Row end: \vrs{value_push} on the vector
and on $b$'s prefix, and distributivity. Exit: $b[0..m] = b$, then
\vrs{normalize}. \dfy{MulLoop}: entry by \dfy{b[0..] == b}; one step unfolds
\dfy{MulAcc} once, using \dfy{b[k..][1..] == b[k + 1..]}; at $k = |b|$ the
invariant says \dfy{r == Times(a, b)}; H1's lemma gives the value conjuncts.
(d) As named in (c), all from L1--L4, L6 and \vrs{vstd::arithmetic::mul}.
(e) No gap.
\end{answer}

% ---------------------------------------------------------------------------
\section*{Problem B2 (paper card)}\label{ans:B2}
\begin{answer}
\textbf{Source:} cheatsheet entry 01.13, lines 458--481 (the pinned range).

\textbf{Task shape:} Isabelle competition problems (miniF2F: 488, split evenly into
validation and test) proved by turning an informal proof, a human's or a
model's, into a formal sketch with open conjectures that automation then
closes.

\textbf{Oracle:} Isabelle itself: an attempt succeeds only when the prover closes
every open conjecture and the filled-in proof checks.

\textbf{What is trusted:} Only the proof checker; the draft and the sketch are not trusted,
and the card says the verifier is the trust anchor.

\textbf{Cheat surface:} A correct formal proof can sit on a wrong draft: of 50 checked model
drafts behind valid proofs, 21 were flawed (13 solved by the prover alone, 8
corrected while formalizing). The prompt demos come from the validation
split.

\textbf{Metric:} Success rate on validation and test: DSP with human drafts 42.6 and
39.3, with Minerva 62B drafts 43.9 and 37.7, against 18.0 and 20.9 for
Sledgehammer with heuristics and 37.3 and 35.2 for Thor with expert
iteration.

\textbf{Budget:} 100 drafts sampled per problem; at 200 attempts model drafts solve
209 problems and human ones 203; compute is not matched across the
baselines.

\textbf{Variance:} Not given in the range: no seeds or intervals; repeated runs would
be needed.

\textbf{Artifact:} The aligned informal proofs are released.
\end{answer}

\end{document}
